% Vorlage fuer Abschlussarbeiten am Fachbereich Mathematik der TU Darmstadt.
% Geeignet fuer Bachelorarbeiten, Masterarbeiten und Diplomarbeiten.

%\documentclass[12pt,a4paper,twoside]{report}
\documentclass[12pt,a4paper,twoside]{article}
% Benutzt man die Option draft, so kann man die Umbrueche ueberpruefen.

% Hier werden alle benoetigten Pakete und Einstellungen geladen. Auch
% hier sind Sie frei diese direkt in der Praeambel zu laden.


\usepackage[ngerman]{babel}

% Kodierung fuer unixoide Systeme und Windowssysteme.

\usepackage[utf8]{inputenc}

% Schriftart Times New Roman.

\usepackage{times}

% Erweitert den Zeichenvorrat, so dass z.B. auch Umlaute im PDF-Dokument
% gefunden werden.

\usepackage[T1]{fontenc}

% Zur Einbindung von Bildern.

\usepackage{graphicx}

% Erweiterte enumerate-Umgebung.

\usepackage{enumerate}

% Verschiedene Pakete, die nuetzlich sind um Mathematik in LaTeX zu setzen.

\usepackage{amsmath, amssymb, amsthm, dsfont}

% Indentation

\usepackage{scrextend}

% Spaces for newcommand

\usepackage{xspace}

% tikz
\usepackage{tikz}
\usetikzlibrary{arrows.meta}% arrow tip library
\usetikzlibrary{bending}% better arrow head for bended lines

% Hyphenation of wordas that already contain a hyphen
% http://tex.stackexchange.com/questions/2706/adequate-hyphenation-of-words-already-containing-a-hyphen
% Example:
% \textsc{Alexandroff}\hyp{}\textsc{Urysohn}\hyp{}Kompaktheit

\usepackage{hyphenat}

% Hyperlinks

%\usepackage{hyperref}
\usepackage[pdfauthor={Fabian Gabel},%
pdftitle={Seminar Lie-Algebren - Wurzelsysteme: einfache Wurzeln, die Weyl-Gruppe und Irreduzibilität},%
%pagebackref=true,%
%pdftex
]{hyperref}

%%%%%%%%%%%%%%%%
% Seitenlayout %
%%%%%%%%%%%%%%%%

% DIV# gibt den Divisor für die Layoutberechnung an.
% Vergrößern des Divisors vergrößert den Textbereich.
% BCOR#cm gibt die Breite des Bundstegs an.
\usepackage[DIV14,BCOR2cm]{typearea}

% Abstand obere Blattkante zur Kopfzeile ist 2.54cm - 15mm
\setlength{\topmargin}{-15mm}

% Keine Einrueckung nach einem Absatz.

\parindent 0pt

% Abstand zwischen zwei Abs\"atzen.

\parskip 12pt

% Zeilenabstand.

\linespread{1.25}

% Inhaltsverzeichnis erstellen.

\usepackage{makeidx}
\makeindex


% Hier werden Makros und Umgebungen eingebunden. Diese werden separat
% in der Datei befehle.tex definiert. Sie sind frei diese Befehle auch
% direkt in der Praeambel zu definieren.
% Makros

\newcommand{\N}{\mathbb{N}} % natuerliche Zahlen
\newcommand{\Z}{\mathbb{Z}} % ganze Zahlen
\newcommand{\Q}{\mathbb{Q}} % rationale Zahlen
\newcommand{\R}{\mathbb{R}} % reelle Zahlen
\newcommand{\co}{\mathrm{co}} % compact open

\DeclareMathOperator{\Ret}{Re} % Realteil
%\newcommand{\PIT}{\textbf{PIT }} % Prime Ideal Theorem
\newcommand{\AC}{\textbf{AC}\xspace} % Prime Ideal Theorem
\newcommand{\UFT}{\textbf{UFT}\xspace} % Prime Ideal Theorem
\newcommand{\PIT}{\textbf{PIT}\xspace}
\newcommand{\ZF}{\textbf{ZF}\xspace}
\newcommand{\ZFC}{\textbf{ZFC}\xspace}
\newcommand{\CC}{\textbf{CC}\xspace}
\newcommand{\CCR}{\textbf{CC}($\R$)\xspace}
\newcommand{\PCC}{\textbf{PCC}\xspace}
\newcommand{\PCCR}{\textbf{PCC}($\R$)\xspace}

\newcommand{\lie}{\textsc{Lie}}

% Umgebungen für Definitionen, Sätze, usw.

%\newtheorem{defn}{Definition}[section]
%\newtheorem{ex}{Beispiel}[chapter]

\theoremstyle{plain}
%\newtheorem{thm}{Satz}[chapter]
\newtheorem{thm}{Satz}[section]
\newtheorem{lem}[thm]{Lemma}
\newtheorem{cor}[thm]{Korollar}
\newtheorem{prop}[thm]{Proposition}

\theoremstyle{definition}
\newtheorem{defn}[thm]{Definition}

\theoremstyle{remark}
\newtheorem*{bem}{Bemerkung}

\def\Satzrefname{Satz}


\begin{document}
% Auf der Titelseite und im Inhaltsverzeichnis sollen keine
% Seitenzahlen dargestellt werden.
\pagestyle{empty}

% Hier wird die Titelseite eingebunden.
\begin{titlepage}
  \begin{center}
    \vspace{1cm}
    \includegraphics[width=0.5\linewidth]{TU_Darmstadt_Logo.pdf}
    \vspace{12pt}
    
    \large{Fachbereich Mathematik}
    \vspace{2cm}
    
    \large{Seminar zu Lie-Algebren}
    \vspace{2cm}

    \huge{Wurzelsysteme: einfache Wurzeln, Weyl-Gruppe und Irreduzibilität}
    
    \vspace*{2cm}    
    
		\large
    Fabian Gabel
    \vspace*{.5cm}

    29.05.2016 \\
    \vspace*{1cm}

    Betreuer: Prof. Dr. rer. nat. Jan-Hendrik Bruinier,  \\
    \hspace{0.6cm}M.Sc. Markus Schwagenscheidt

    \vspace*{.5cm}

    %Zweiter Gutachter: Name des zweiten Gutachters\\[2ex]
    \tiny{Version vom \today}
  \end{center}
\end{titlepage}
\vspace*{\fill}


% Inhaltsverzeichnis erstellen.
\tableofcontents

% Ab sofort werden Seitenzahlen in der Kopfzeile angezeigt.
\pagestyle{headings}

\section*{Einleitung}\index{Einleitung}
\addcontentsline{toc}{section}{Einleitung}\index{Einleitung}

Basierend auf \cite[S.49-55]{humphreys1972introduction} soll in dieser Ausarbeitung auf . . . eingegangen werden. 

\section{Grundlagen zu Wurzelsystemen}




\section{Einfache Wurzeln}
\label{sec:einfach}

In diesem Abschnitt sollen einige Eigenschaften einfacher Wurzeln bewiesen werden.  
Im Folgenden bezeichne $\Delta$ ein fest gewähltes Fundamentalsystem des Wurzelsystems $\Phi$ im \euklid ischen Vektorraum $E$.

\begin{lem}
  \label{lem:independentSet}
  Sei $M \subseteq E$ eine Teilmenge von Vektoren und $P_\gamma$ eine Hyperebene.
  Es gelte $(\alpha, \gamma) > 0$ für alle $\alpha \in M$ und zudem $(\alpha, \beta) \leq 0$ für alle paarweise verschiedenen $\alpha, \beta \in M$.
  Dann ist $M$ linear unabhängig.
\end{lem}

\begin{proof}
  Wir betrachten die Gleichung $\sum_{\alpha \in M} k_\alpha \alpha = 0$ und definieren die disjunkten Mengen $A := \{\alpha \in M \mid k_\alpha < 0\}$ und  $B := \{\beta \in M \mid k_\beta > 0\}$. 
  Damit gilt 
  \begin{displaymath}
    \sum_{\alpha \in A} k_\alpha \alpha = \sum_{\beta \in B} k_\beta \beta \tag{$\ast$}.
  \end{displaymath}
  Daraus folgt mit $\varepsilon := \sum_{\alpha \in A} k_\alpha \alpha$ dann
  \begin{displaymath}
    0 
    \leq (\varepsilon, \varepsilon) 
    = (\;\sum_{\alpha \in A} k_\alpha \alpha,\sum_{\beta \in B} k_\beta \beta\;) 
    = \sum_{\alpha \in A, \beta \in B} k_\alpha k_\beta \; (\alpha, \beta) 
    \leq 0,
  \end{displaymath}
  aufgrund der Eigenschaften des Skalarprodukts und der Voraussetzung $(\alpha, \beta) \leq 0$ für alle paarweise verschiedenen $\alpha, \beta$.
  Dies impliziert $\varepsilon = 0$, womit auch 
  \begin{displaymath}
    0 = (\varepsilon, \gamma) = \sum_{\alpha \in A} k_\alpha (\alpha, \gamma)
  \end{displaymath}
  folgt. Da voraussetzungsgemäß $(\alpha, \gamma) > 0$ gilt, müssen alle $k_\alpha$ mit $\alpha \in A$ bereits $0$ sein.
  Analog folgert man unter Verwendung der Identität ($\ast$), dass alle $k_\beta$ mit $\beta \in B$ gleich $0$ sein müssen.
\end{proof}

Das folgende Lemma demonstriert, wie sich aus einer positiven Wurzel neue positive Wurzeln gewinnen lassen.

\begin{lem}
  \label{lem:posRootNewRoot}
  Ist $\alpha \in \Phi$ eine positive aber nicht einfache Wurzel, so existiert ein $\beta \in \Delta$, sodass die Differenz $\alpha - \beta$ eine positive Wurzel ist.
\end{lem}

\begin{proof}
  Wir wollen zeigen, dass ein $\beta \in \Delta$ existiert mit $(\alpha, \beta) > 0$.
  Angenommen, für alle $\beta \in \Delta$ gelte $(\alpha, \beta) \leq 0$.
  Zudem gilt für alle $\beta, \beta' \in \Delta$ immer $(\beta, \beta') \leq 0$, denn sonst wäre nach Lemma \ref{lem:sumDiffRoot} auch $\beta - \beta'$ eine Wurzel, was jedoch \hyperref[it:B2]{(B2)} widerspricht.
  Damit sind die Voraussetzungen für Lemma \ref{lem:independentSet} erfüllt und es ist $\Delta \cup \{\alpha\}$ eine linear unabhängige Menge.
  Dies widerspricht jedoch der Voraussetzung, dass $\Delta$ ein Fundamentalsystem ist.

  Also muss ein $\beta \in \Delta$ existieren mit $(\alpha, \beta) > 0$.
  Zudem kann $\beta$ nicht proportional zu $\alpha$ sein, da sonst $\alpha$ nicht positiv sein könnte.
  Unter Verwendung von Lemma \ref{lem:sumDiffRoot} folgt somit, dass $\alpha - \beta$ eine Wurzel ist.
  
  Bezüglich $\Delta$ lässt sich $\alpha$ auch schreiben als $\alpha = \sum_{\gamma \in \Delta} k_\gamma \gamma$, wobei aufgrund der Positivität von $\alpha$ alle $k_\gamma \geq 0$ sind und darüber hinaus auch mindestens ein $\gamma \neq \beta$ mit  $k_\gamma > 0$ existiert.
  Ansonsten wäre nämlich $\alpha$ proportional zu $\beta$. 
  Damit würde jedoch nach \hyperref[it:R2]{(R2)} entweder $\alpha = \beta$ gelten, was der Voraussetzung, dass $\alpha$ nicht einfach ist, widerspricht, oder es würde $\alpha = -\beta$ gelten, was der vorausgesetzten Positivität von $\alpha$ widerspricht.

  Die Eigenschaft \hyperref[it:B2]{(B2)} erzwingt nun die Positivität aller restlichen Linearfaktoren und folglich ist auch $\alpha - \beta$ eine positive Wurzel.
\end{proof}
 
\begin{cor}
  Jedes $\beta \in \Phi^+$ lässt sich als Linearkombination $\alpha_1 + \dots + \alpha_k$ mit $\alpha_i \in \Delta$ so schreiben, dass für alle $i \in \{1,\dots,k\}$ die Partialsumme $\alpha_1 + \dots + \alpha_i$ eine Wurzel ist.
\end{cor}

\begin{proof}
  Wir führen einen Induktionsbeweis und induzieren über die Höhe $\hgt(\beta)$.
  Sei $\hgt(\beta) = 1$, dann ist $\beta \in \Delta$ und die Behauptung ist erfüllt.

  Angenommen, die Behauptung gelte für ein $k-1 \in \N$.
  Sei nun $\hgt(\beta) = k$ und damit $\beta$ positiv aber nicht einfach.
  Nach \hyperref[it:B2]{(B2)} lässt sich dann 
  \begin{displaymath}
    \beta = \alpha_1 + \dots + \alpha_k \tag{$\ast$}
  \end{displaymath}
  schreiben, wobei die $\alpha_i$ nicht zwingend paarweise verschieden sind.
  Mit Lemma \ref{lem:posRootNewRoot} folgt dann, dass ein $\alpha \in \Delta$ existiert, sodass $\beta - \alpha$ eine positive Wurzel ist.
  Aufgrund der Eindeutigkeit der Darstellung ($\ast$) bis auf Reihenfolge der Summanden, lässt sich ohne Beschränkung der Allgemeinheit $\alpha = \alpha_k$ annehmen.

  Es gilt $\hgt(\beta - \alpha_{k}) = k - 1$ und nach Induktionsvoraussetzung folgt $\alpha_1 + \dots + \alpha_{k-1} \in \Phi^+$ für alle $i \in \{1,\dots,k-1\}$.
  Der Fall $i = k$ folgt jedoch bereits mit ($\ast$), womit die Behauptung für $k$ vollständig bewiesen ist.
\end{proof}

\begin{lem}
  \label{lem:permute}
  Sei $\alpha \in \Delta$. 
  Dann permutiert die Spiegelung $\sigma_\alpha$ alle von $\alpha$ verschiedenen positiven Wurzeln, es gilt also $\sigma_\alpha(\Phi^+ \setminus \{\alpha\}) = \Phi^+ \setminus \{\alpha\}$.
\end{lem}

\begin{proof}
  Sei $\beta \in \Phi^+ \setminus \{\alpha\}$. 
  Dann existiert eine Darstellung der Form $\beta = \sum_{\gamma \in \Delta} k_\gamma \gamma$ mit positiven Koeffizienten $k_\gamma$.
  Es kann $\beta$ nicht proportional zu $\alpha$ sein, da sonst nach \hyperref[it:R2]{(R2)} und unter der Voraussetzung $\beta \neq \alpha$ ein Widerspruch zur Voraussetzung entsteht, dass $\beta$ positiv ist.
  Daher existiert ein $\gamma \neq \beta$ mit $k_\gamma > 0$.

  Aufgrund der Linearität von $\sigma_\alpha$ besitzt dann auch $\sigma_\alpha(\beta)$ mindestens einen positiven Linearfaktor und damit ist nach \hyperref[it:B2]{(B2)} auch $\sigma_\alpha(\beta)$ positiv.
  Des Weiteren ist $\sigma_\alpha(\beta)$ nicht proportional zu $\alpha$, da sonst $\beta = -\alpha$ gelten würde, was im vorangehenden Absatz bereits ausgeschlossen wurde.
  Hieraus folgt die Behauptung.
\end{proof}

\begin{cor}
  \label{cor:sigmaDelta}
  Sei $\delta := \tfrac{1}{2} \sum_{\beta \in \Phi^+} \beta$.
  Dann gilt $\sigma_\alpha(\delta) = \delta - \alpha$ für alle $\alpha \in \Delta$.
\end{cor}

\begin{proof}
  Es gilt mit Lemma \ref{lem:permute}
  \begin{displaymath}
    \frac{1}{2} \sum_{\substack{\beta \in \Phi^+\\ \beta \neq \alpha}} \sigma_\alpha(\beta) + \frac{1}{2}\alpha 
  = \frac{1}{2} \sum_{\substack{\beta \in \Phi^+\\ \beta \neq \alpha}} \beta + \frac{1}{2}\alpha
  = \frac{1}{2} \sum_{\beta \in \Phi^+} \beta
  = \delta.
  \end{displaymath}
  Daraus folgt 
  \begin{align*}
    \sigma_\alpha(\delta) 
    &= \sigma_\alpha(\frac{1}{2} \sum_{\substack{\beta \in \Phi^+\\ \beta \neq \alpha}} \sigma_\alpha(\beta) + \frac{1}{2}\alpha ) \\
    &= \frac{1}{2} \sum_{\substack{\beta \in \Phi^+\\ \beta \neq \alpha}} \sigma_\alpha(\beta) + \frac{1}{2} \sigma_\alpha(\alpha) \\
    &= \frac{1}{2} \sum_{\substack{\beta \in \Phi^+\\ \beta \neq \alpha}} \sigma_\alpha(\beta) + \frac{1}{2} \sigma_\alpha(\alpha) + \frac{1}{2} \alpha - \frac{1}{2} \alpha \\
    &= \frac{1}{2} \sum_{\substack{\beta \in \Phi^+\\ \beta \neq \alpha}} \beta + \frac{1}{2} \alpha - \alpha \\
    &= \frac{1}{2} \sum_{\beta \in \Phi^+} \beta - \alpha \\
    &= \delta - \alpha. \qedhere
  \end{align*}
\end{proof}

Nach Definition \ref{def:weylgroup} wird die \weyl\hyp{}Gruppe $\W$ von den Spiegelungen an den zu Wurzeln orthogonalen Hyperebenen erzeugt.
Das folgende Lemma liefert ein Kriterium dafür, wann man entsprechende Darstellungen von Elementen aus $\W$ vereinfachen kann.

\begin{lem}
  \label{lem:simplify}
  Seien $\alpha_1,\dots,\alpha_t$ nicht notwendig verschiedene einfache Wurzeln. Es sei $\sigma_i := \sigma_{\alpha_i}$.
  Ist $\sigma_1 \dots \sigma_{t-1}(\alpha_t)$ eine negative Wurzel, dann existiert ein Index $1 \leq s < t$, sodass $\sigma_1 \dots \sigma_t = \sigma_1 \dots \sigma_{s-1} \sigma_{s+1} \dots \sigma_{t-1}$ gilt.
\end{lem}

\begin{proof}
  Sei $\beta_i := \sigma_{i + 1} \dots \sigma_{t-1}(\alpha_t)$, für alle $0 \leq i \leq t - 2$ und $\beta_{t - 1} := \alpha_t$.
  Nach Voraussetzung ist $\beta_0$ negativ und $\beta_{t - 1}$ als einfache Wurzel positiv.
  Es existiert also aufgrund der Endlichkeit der Folge der $\beta_i$ ein minimaler Index $s$, sodass jeweils $\beta_s$ positiv und $\sigma_s(\beta_s) = \beta_{s - 1}$ negativ sind.
  Wiederum impliziert \ref{lem:permute} die Gleichung $\beta_s = \alpha_s$, da sonst andernfalls $\sigma_s\beta_s$ positiv sein müsste.
  Mit Lemma \ref{lem:conjReflection} folgt sodann
  \begin{displaymath}
    \sigma_s 
    = \sigma_{\alpha_s} 
    = \sigma_{\sigma_{s+1} \dots \sigma_{t - 1}(\alpha_t)}
    = (\sigma_{s+1} \dots \sigma_{t - 1}) \sigma_{t} (\sigma_{s+1} \dots \sigma_{t - 1})^{-1}.
  \end{displaymath}
  Dies impliziert wiederum
  \begin{align*}
    \sigma_1 \dots \sigma_t 
    &= \sigma_1 \dots \sigma_{s - 1} \sigma_s \sigma_{s+1} \dots \sigma_t \\
    &= \sigma_1 \dots \sigma_{s - 1} (\sigma_{s+1} \dots \sigma_{t - 1}) \sigma_{t} (\sigma_{s+1} \dots \sigma_{t - 1})^{-1} \sigma_{s+1} \dots \sigma_{t - 1} \sigma_t\\
    &= \sigma_1 \dots \sigma_{s - 1} (\sigma_{s+1} \dots \sigma_{t - 1}) \sigma_t  \sigma_t\\
    &= \sigma_1 \dots \sigma_{s - 1} \sigma_{s+1} \dots \sigma_{t - 1}. \qedhere
  \end{align*}
\end{proof}

\begin{cor}
  \label{cor:minimalReflection}
  Sei $\sigma = \sigma_1 \dots \sigma_t \in \W$, mit einfachen Spiegelungen $\sigma_i$, wobei $t$ minimal gewählt sei. Dann ist $\sigma(\alpha_t)$ eine negative Wurzel.
\end{cor}

\begin{proof}
  Nach Voraussetzung lässt sich der Ausdruck für $\sigma$ nicht weiter verkürzen.
  Mit Lemma \ref{lem:simplify} folgt, dass $\sigma_1 \dots \sigma_{t-1}(\alpha_t)$ positiv ist.
  Da $\sigma_t$ eine Spiegelung ist, gilt $\alpha_t = \sigma_t(-\alpha_t)$.
  Damit folgt dann, dass auch $\sigma_1 \dots \sigma_{t-1}(\sigma_t(-\alpha_t)) = \sigma_1 \dots \sigma_t(-\alpha_t)$ positiv ist.
  Die Linearität von Spiegelungen impliziert dann die Negativität von $\sigma = \sigma_1 \dots \sigma_t(\alpha_t)$.
\end{proof}

\section{Die Weyl-Gruppe}
\label{sec:weylgroup}

Nach der in den bisherigen Abschnitten geleisteten Vorarbeit, wenden wir uns in diesem Abschnitt nun wieder der Gruppenoperation der \weyl\hyp{}Gruppe auf der Menge der \weyl\hyp{}Kammern beziehungsweise auf der Menge der Fundamentalsysteme eines Wurzelsystems $\Phi$ aus Proposition \ref{prop:groupOp} zu.

Wir beginnen mit einer für beliebige Gruppenoperationen definierten Eigenschaft.

\begin{defn}
  Eine Gruppe $G$ operiere auf einer Menge $X$.
  Die Gruppenoperation $\circ \colon G \times X \to X$ heißt \emph{scharf transitiv}, wenn für alle $x,y \in X$ genau ein $g \in G$ existiert, sodass $g \circ x = y$. 
\end{defn}

Ziel dieses Abschnittes ist es nun zu beweisen, dass die durch $\W$ auf der Menge der Fundamentalsysteme induzierte Gruppenoperation scharf transitiv ist.
Dies wird unter anderem auch durch den nachfolgenden Satz bewiesen.
Ein Lemma stellt zuvor noch Sicher, dass es sich bei $\W$ um eine endliche Gruppe handelt.

\begin{lem}
  \label{lem:weylFinite}
  Es sei $\Phi$ ein Wurzelsystem in $E$. Dann ist die entsprechende \weyl\hyp{}Gruppe endlich. 
\end{lem}

\begin{proof}
  Aus \hyperref[it:R3]{(R3)} folgt, dass $\W$ auf der Menge aller Wurzeln operiert, welche nach \hyperref[it:R1]{(R1)} endlich ist.
  Es existiert also ein Gruppenhomomorphismus in die symmetrische Gruppe $\Sym(\Phi)$.

  Ist dieser Homomorphismus injektiv, also die Gruppenoperation treu, so lässt sich $\W$ als eine Untergruppe von $\Sym(\Phi)$ auffassen und ist damit endlich.
  Sei $\sigma \in \W$ im Kern dieses Homomorphismus.
  Dann fixiert $\sigma$ alle Wurzeln aus $\Phi$.
  Da nach \hyperref[it:R1]{(R1)} die Wurzeln den Vektorraum $E$ aufspannen ist $\sigma$ bereits eindeutig festgelegt und es folgt $\sigma = \mathds{1}$.
  Also ist die Gruppenoperation treu und $\W$ damit endlich.
\end{proof}


\begin{thm}
  \label{thm:simplyTransitive}
  Es sei $\Delta$ ein Fundamentalsystem des Wurzelsystems $\Phi$ in $E$.
  Dann gelten die folgenden Aussagen:
  \begin{enumerate}[(a)]

    \item Wenn $\gamma \in E$ regulär ist, dann existiert ein $\sigma \in \W$, sodass $(\sigma(\gamma), \alpha) > 0$ für alle $\alpha \in \Delta$. 
      Es operiert $\W$ also transitiv auf der Menge der \weyl\hyp{}Kammern.

    \item Wenn $\Delta'$ ein weiteres Fundamentalsystem von $\Phi$ ist, dann existiert ein $\sigma \in \W$, sodass $\sigma(\Delta') = \Delta$.
      Es operiert $\W$ also transitiv auf der Menge der Fundamentalsysteme.

    \item Für eine beliebige Wurzel $\alpha \in \Phi$ existiert ein $\sigma \in \W$, sodass $\sigma(\alpha) \in \Delta$.
      Jede Wurzel liegt also in der $\W$\hyp{}Bahn einer einfachen Wurzel.
    \item Die \weyl\hyp{}Gruppe wird erzeugt von den Spiegelungen $\sigma_\alpha$ für $\alpha \in \Delta$.

    \item Ist $\sigma \in \W$ und gilt $\sigma(\Delta) = \Delta$, so folgt $\sigma = \mathds{1}$.
      Also operiert $\W$ einfach transitiv auf der Menge der Fundamentalsysteme.
  \end{enumerate}
\end{thm}

\begin{proof}
  Wir führen die Beweise für (a) bis (c) zunächst für die von den einfachen Spiegelungen erzeugte Untergruppe $\W'$ von $\W$.
  Nach dem Beweis von (d) folgt $\W' = \W$.

  (a):
  Es sei $\gamma \in E$ regulär und $\delta := \tfrac{1}{2} \sum_{\alpha \in \Phi^+} \alpha$.
  Da nach Lemma \ref{lem:weylFinite} die \weyl\hyp{}Gruppe und damit auch die Untergruppe $\W'$ endlich sind, existiert ein $\sigma \in \W'$, sodass $(\sigma(\gamma), \delta)$ maximal wird.
  Für $\alpha \in \Delta$ gilt auch $\sigma_\alpha\sigma \in \W'$ aufgrund der Untergruppeneigenschaft.
  Daraus folgt
  \begin{align*}
    (\sigma(\gamma), \delta) 
    &\geq (\sigma_\alpha \sigma(\gamma), \delta) &&\text{aufgrund der Maximalität} \\
    &= (\sigma(\gamma), \sigma_\alpha(\delta)) &&\text{da $\sigma_\alpha$ eine Isometrie und zudem selbstinvers ist} \\
    &= (\sigma(\gamma), \delta - \alpha) &&\text{nach Korollar \ref{cor:sigmaDelta}} \\
    &= (\sigma(\gamma), \delta) - (\sigma(\gamma), \alpha) &&\text{aufgrund der Bilinearität des Skalarproduktes.}
  \end{align*}
  Also muss $(\sigma(\gamma), \alpha) \geq 0$ für alle $\alpha \in \Delta$ gelten, um nicht der Maximalität von $\sigma$ zu widersprechen.

  Nach Voraussetzung ist $\gamma$ regulär.
  Daher kann für kein $\alpha \in \Delta$ die Identität $(\sigma(\gamma),\alpha) = 0$ gelten, da sonst $\gamma$ in $P_{\sigma^{-1}(\alpha)}$ liegt, was der vorausgesetzten Regularität widerspricht.
  Somit gilt für alle $\alpha \in \Delta$ die strikte Ungleichung $(\sigma(\gamma), \alpha) > 0$.
  Das bedeutet jedoch gerade, dass $\sigma(\gamma)$ in der Fundamentalkammer $\mathfrak{C}(\Delta)$ liegt.

  Also bildet entsprechend Proposition \ref{prop:groupOp} die Spiegelung $\sigma$ die \weyl\hyp{}Kammer $\mathfrak{C}(\gamma)$ auf $\mathfrak{C}(\Delta)$ ab. 
  Somit gilt, dass alle \weyl\hyp{}Kammern mit der Fundamentalkammer verbunden sind, die Menge der \weyl\hyp{}Kammern ist also eine $\W'$\hyp{}Bahn und die Gruppenoperation damit transitiv.


  (b):
  Da nach dem vorangehenden Beweis $\W'$ die \weyl\hyp{}Kammern transitiv permutiert, gilt dies nach Proposition \ref{prop:groupOp} auch für die korrespondierende Gruppenoperation auf Fundamentalsystemen.

  (c):
  Da nach (b) die Untergruppe $\W'$ transitiv auf der Menge der Fundamentalsysteme operiert, genügt es nachzuweisen, dass jede Wurzel in einem Fundamentalsystem enthalten ist.
  UNVOLLSTÄNDIG

  (d):
  Da nach Konstruktion $\W'$ eine Untergruppe von $\W$ ist, reicht es zu zeigen, dass jede Spiegelung $\sigma_\alpha$ mit $\alpha \in \Phi$ ein Element von $\W'$ ist.
  Nach (c) existiert ein $\sigma \in \W'$, sodass $\beta := \sigma(\alpha) \in \Delta$.
  Damit gilt insbesondere $\sigma_\beta \in \W'$.
  Aufgrund Lemma \ref{lem:conjReflection} ist $\sigma_{\sigma(\alpha)} = \sigma \sigma_\alpha \sigma^{-1}$, also folgt mit der Untergruppeneigenschaft von $\W'$ sodann $\sigma_\alpha = \sigma^{-1} \sigma_\beta \sigma \in \W'$.
  Folglich gilt bereits $\W' = \W$ und alle bereits für $\W'$ bewiesenen Aussagen gelten auch für $\W$.

  (e):
  Sei $\sigma \in \W$ mit $\sigma(\Delta) = \Delta$.
  Angenommen, $\sigma \neq \mathds{1}$.
  Nach (d) lässt sich $\sigma$ als Produkt einfacher Spiegelungen schreiben, es gilt also $\sigma = \sigma_1\dots\sigma_t$ und $\sigma_1 \dots \sigma_t(\alpha_t)$. 
  Unter Berücksichtigung von Lemma \ref{lem:simplify} kann man Annehmen, dass die multiplikative Darstellung nicht mehr verkürzt werden kann.
  Dann ist nach Korollar \ref{cor:minimalReflection} das Bild $\sigma(\alpha_t)$ negativ, also gilt insbesondere $\sigma(\alpha_t) \not\in \Delta$ im Widerspruch zur Voraussetzung.
  Die Annahme muss also verworfen werden und es folgt $\sigma = \mathds{1}$.
\end{proof}

ÜBERLEITUNG

\begin{defn}
  Sei $\Phi$ ein Wurzelsystem mit Fundamentalsystem $\Delta$ und \weyl\hyp{}Gruppe $\W$. 
  Des Weiteren sei $\sigma \in \W$, mit einer Darstellung $\sigma = \sigma_{\alpha_1}\dots\sigma_{\alpha_t}$ ($\ast$), wobei $t$ minimal gewählt sei und $\alpha_i \in \Delta$. 
  Dann nennt man den Ausdruck ($\ast$) auch \emph{reduziert} und definiert die \emph{Länge von $\sigma$ bezüglich} $\Delta$ als $\ell(\sigma) := t$.
\end{defn}

\begin{lem}
  \label{lem:lengthAndNegativeRoots}
  Sei $\Phi$ ein Wurzelsystem in $E$ und $\W$ die zugehörige \weyl\hyp{}Gruppe.
  Für $\sigma \in \W$ bezeichne $n(\sigma)$ die Anzahl der positiven Wurzeln $\alpha$ mit negativem Bild $\sigma(\alpha)$.
  Dann gelten die folgenden Aussagen:
  \begin{enumerate}[(1)]
    \item Sei $\alpha \in \Delta$. Falls $\sigma(\alpha)$ negativ ist, so gilt $n(\sigma\sigma_\alpha(\alpha)) = n(\sigma) - 1$.
    \item Es gilt $\ell(\sigma) = n(\sigma)$.
  \end{enumerate}
\end{lem}

\begin{proof}
  (1):
  Sei $\sigma(\alpha)$ negativ.
  Dann folgt mi Lemma \ref{lem:permute}, dass $\sigma_\alpha$ alle von $\alpha$ verschiedenen positiven Wurzeln permutiert.
  Es enthält somit $\sigma(\Phi^+ \setminus \{\alpha\})$ genauso viele negative Wurzeln wie $\sigma(\sigma_\alpha(\Phi^+ \setminus \{\alpha\}))$.
  Andererseits ist jedoch $\sigma \sigma_\alpha (\alpha) = -\sigma(\alpha)$ positiv.
  Daraus können wir $n(\sigma\sigma_\alpha) = n(\sigma) - 1$ folgern.

  (2):
  Wir führen den Beweis mittels Induktion über $\ell(\sigma)$.

  Ist $\ell(\sigma)=0$, so gilt $\sigma = \mathds{1}$, was wiederum $n(\sigma)=0$ zur Folge hat.

  Angenommen, die Behauptung gelte für alle $\tau \in \W$ mit $\ell(\tau) < \ell(\sigma)$.
  Wir schreiben nun $\sigma$ in der reduzierten Form $\sigma = \sigma_{\alpha_1} \dots \sigma_{\alpha_t}$ und definieren $\alpha := \alpha_t$.
  Mit Korollar \ref{cor:minimalReflection} folgt damit die Negativität von $\sigma(\alpha)$.
  Aufgrund von (1) können wir einerseits $n(\sigma\sigma_\alpha) = n(\sigma) - 1$ folgern.
  
  Andererseits gilt $\ell(\sigma\sigma_\alpha) = \ell(\sigma_{\alpha_1} \dots \sigma_{\alpha_{t -1}}) = \ell(\sigma) - 1 < \ell(\sigma)$.
  Mit der Induktionsvoraussetzung folgt somit $\ell(\sigma\sigma_\alpha) = n(\sigma\sigma_\alpha)$, also $\ell(\sigma) - 1 = n(\sigma) - 1$, woraus sich die Behauptung ergibt.
\end{proof}

Nun betrachten wir die einfach transitive Operation der \weyl\hyp{}Gruppe auf den \weyl\hyp{}Kammern aus Satz \ref{thm:simplyTransitive}.
Es soll gezeigt werden, dass der in der \textsc{euklid}ischen Topologie gebildete Abschluss $\overline{\mathfrak{C}(\Delta)}$ der Fundamentalkammer bezüglich des Fundamentalsystems $\Delta$ ein \emph{Fundamentalbereich} der Gruppenoperation von $\W$ auf $E$ ist.
Dies bedeutet, dass für alle Vektoren $x \in E$ genau ein $y \in \overline{\mathfrak{C}(\Delta)}$ existiert mit $x \in \W y = \{\sigma(y) \mid \sigma \in \W\}$.

\begin{lem}
  \label{lem:fundamentalDomain}
  Sei $\Delta$ ein Fundamentalsystem eines Wurzelsystems $\Phi$ in $E$ und $\W$ die entsprechende \weyl\hyp{}Gruppe.
  Dann gelten folgende Aussagen:
  \begin{enumerate}[(1)]
    \item Jeder Punkt aus $E$ liegt in der $\W$\hyp{}Bahn eines Elementes aus  $\overline{\mathfrak{C}(\Delta)}$.
    \item Falls für $\lambda, \mu \in \overline{\mathfrak{C}(\Delta)}$ ein $\sigma \in \W$ existiert mit $\sigma \lambda = \mu$, dann ist $\sigma$ das Produkt einfacher Spiegelungen, die $\lambda$ fixieren und es gilt speziell $\lambda = \mu$.
  \end{enumerate}
  Insbesondere ist $\overline{\mathfrak{C}(\Delta)}$ ein Fundamentalbereich der Gruppenoperation von $\W$ auf $E$.
\end{lem}

\begin{proof}
  (1):
  Wir setzen die in Abschnitt \ref{sec:grundlagen} auf der Menge der Wurzeln $\Phi$ definierte Halbordnungsrelation $\succeq$ nun auf ganz $E$ fort, indem wir genau dann $\beta \preceq \alpha$ setzen, wenn $\alpha - \beta$ eine $\R$-Linearkombination mit sämtlich positiven Koeffizienten ist oder $\alpha = \beta$ gilt.

  Sei $\lambda \in E$ gegeben. 
  Wir betrachten als nächstes die Teilmenge $M$ aller $\W$\hyp{}Bahnelemente $\mu \in \W\lambda$, für die $\mu \succeq \lambda$ gilt.
  Da $M$ endlich und aufgrund von $\lambda \in M$ nichtleer ist, existieren bezüglich der Halbordnung $\succeq$ maximale Elemente.

  Sei $\mu = \sigma\lambda$ mit $\sigma \in \W$ ein solches maximales Element.
  Wir möchten nun zeigen, dass $\mu \in \overline{\mathfrak{C}(\Delta)}$ gilt.
  Sei dazu $\alpha \in \Delta$.
  Es ist $\sigma_\alpha(\mu) = \mu - \langle \mu, \alpha \rangle \alpha$.
  Da $\sigma_\alpha(\mu) = \sigma_\alpha \sigma \lambda \in \W\lambda$ muss, um nicht der Maximalität von $\mu$ zu widersprechen $\langle \mu, \alpha \rangle \geq 0$ gelten.
  Damit liegt jedoch $\mu$ in $\overline{\mathfrak{C}(\Delta)}$.

  (2):
  Wir führen einen Induktionsbeweis über die Länge $\ell(\sigma)$.

  Ist $\ell(\sigma) = 0$ so folgt $\sigma = \mathds{1}$ und damit sofort $\lambda = \mu$.

  Sei nun $\ell(\sigma) > 0$ und die Aussage für alle $\tau \in \W$ mit $\ell(\tau) < \ell(\sigma)$ bereits bewiesen.
  Nach Lemma \ref{lem:lengthAndNegativeRoots} bildet eine positive Wurzel auf eine negative Wurzel ab.
  Insbesondere existiert also eine einfache Wurzel $\alpha \in \Delta$ mit negativem Bild $\sigma(\alpha)$.
  Für alle $\xi \in \mathfrak{C}(\Delta)$ gilt $(\xi, \sigma(\alpha)) = \sum_{\beta \in \Delta} k_\beta (\xi,\beta) < 0$. 
  Denn alle $k_\beta$ sind nach Voraussetzung negativ sind zudem gilt $(\xi,\beta) > 0$, weil $\xi$ als Element der Fundamentalkammer mit allen einfachen Wurzeln einen spitzen Winkel einschließt.
  Aufgrund der Stetigkeit des Skalarproduktes gilt somit auch $0 \geq (\mu, \sigma(\alpha))$.
  
  Andererseits folgt mit derselben Argumentation angewendet auf $\lambda$ die Identität
  \begin{displaymath}
    0 \geq (\mu, \sigma(\alpha)) = (\sigma^{-1}(\mu), \alpha) = (\lambda, \alpha) \geq 0,
  \end{displaymath}
  weil Spiegelungen Isometrien sind.
  Hieraus folgt zunächst $(\lambda, \alpha) = 0$, also gilt $\lambda \in P_\alpha$.
  Dies wiederum impliziert $\sigma_\alpha(\lambda) = \lambda$ und weiter noch $\sigma \sigma_\alpha(\lambda) = \mu$.
  Somit erfüllt $\sigma \sigma_\alpha$ die Voraussetzungen des zu beweisenden Lemmas. 

  Wir zeigen nun, dass auf $\sigma \sigma_\alpha$ auch die Induktionsannahme zutrifft.
  Nach Lemma \ref{lem:lengthAndNegativeRoots}(1) gilt $n(\sigma\sigma_\alpha) = n(\sigma) - 1$, woraus schließlich mit \ref{lem:lengthAndNegativeRoots}(2) die Gleichheit $\ell(\sigma\sigma_\alpha) = \ell(\sigma) - 1$ folgt.
  Es gilt also $\ell(\sigma\sigma_\alpha) < \ell(\sigma)$ und darum sich aufgrund der Induktionsannahme $\lambda = \mu$ folgern.
\end{proof}
 
Eine Verschärfung von Lemma \ref{lem:fundamentalDomain} findet sich in \cite[S.22]{humphreys1992reflection}.

\section{Irreduzible Wurzelsysteme}
\label{sec:irreduc}

In diesem Abschnitt beschäftigen wir uns mit den Eigenschaften der kleinsten ``Bausteine'' aus denen Wurzelsysteme zusammengesetzt sind.
Es sei wieder $\Phi$ ein fest gewähltes Wurzelsystem über dem \euklid ischen Vektorraum $E$ mit Fundamentalsystem $\Delta$ und zugehöriger \weyl\hyp{}Gruppe $\W$.

\begin{defn}
  \label{def:irreducibleRoot}
  Ein Wurzelsystem $\Phi$ heißt \emph{reduzibel}, wenn eine Zerlegung $\Phi = \Phi_1 \dot{\cup} \Phi_2$ in zwei nichtleere disjunkte Mengen existiert, sodass die Elemente aus unterschiedlichen Mengen paarweise orthogonal sind.
  Ein Wurzelsystem, welches sich nicht orthogonal zerlegen lässt bezeichnet man als \emph{irreduzibel}.
  Entsprechend bezeichnet man auch ein Fundamentalsystem eines Wurzelsystems als \emph{irreduzibel}, falls keine disjunkte Zerlegung im obigen Sinne existiert.
\end{defn}

\begin{bem}
  Sind $(E,\Phi)$ und $(F,\Psi)$ zwei Wurzelsysteme, so ist auch die äußere direkte Summe $E \oplus F$ auf natürliche Weise ein \euklid ischer Vektorraum.
  Ebenso ist die Menge $(R \times \{0_F\}) \cup (\{0_E\} \times S)$ ein Wurzelsystem dieses Produktraumes, welches ebenso als \emph{direkte Summe} der Wurzelsysteme bezeichnet wird.
  Ein Beweis hierzu findet sich in \cite[S.199]{hall2015lie}.
  Ist nun im Sinne von Definition \ref{def:irreducibleRoot} ein Wurzelsystem reduzibel, so liefert die orthogonale Zerlegung des Wurzelsystems eine orthogonale Zerlegung $E = E_1 \oplus E_2$ des Vektorraumes.
  Damit ist ein Wurzelsystem genau dann irreduzibel, wenn es nicht als direkte Summe von zwei Wurzelsystemen geschrieben werden kann.
\end{bem}

\begin{lem}
  \label{lem:irreducibleRoot}
  Das Wurzelsystem $\Phi$ ist genau dann irreduzibel, wenn $\Delta$ irreduzibel ist.
\end{lem}

\begin{proof}
  ``$\Rightarrow$'':
  Angenommen, es existiere eine Zerlegung $\Phi = \Phi_1 \cup \Phi_2$ mit $(\Phi_1, \Phi_2)$ = 0.
  Falls für ein $i \in \{1,2\}$ die Relation $\Delta \not\subseteq \Phi_i$ gilt, ist $\Delta = (\Phi_1 \cap \Delta) \cup (\Phi_2 \cap \Delta)$ eine entsprechende Zerlegung des Fundamentalsystems.
  Andernfalls ist ohne Beschränkung der Allgemeinheit $\Delta \subseteq \Phi_1$ und damit $(\Delta, \Phi_2) = 0$.
  Da nach \hyperref[it:B1]{(B1)} das Fundamentalsystem eine Vektorraumbasis ist, folgt mit der Linearität des Skalarproduktes sogleich $(E, \Phi_2) = 0$ und da Skalarprodukte nicht ausgeartet sind letztlich $\Phi_2 = {0}$ im Widerspruch zu \hyperref[it:B1]{(B1)}.

  ``$\Leftarrow$'':
  Angenommen, $\Phi$ sei irreduzibel und es existiere eine Zerlegung $\Delta = \Delta_1 \cup \Delta_2$ in disjunkte nichtleere Teilmengen $\Delta_1$ und $\Delta_2$ mit $(\Delta_1, \Delta_2) = 0$.
  Wir definieren für alle $i \in \{1,2\}$ die Mengen $\Phi_i$ aller Wurzeln, die in der $\W$\hyp{}Bahn einer Wurzel aus $\Delta_i$ liegen. 
  Nach Satz \ref{thm:simplyTransitive}(c) liegt jede Wurzel in der $\W$\hyp{}Bahn einer einfachen Wurzel, womit $\Phi = \Phi_1 \cup \Phi_2$ gilt.

  Sei $\gamma$ eine Wurzel in $\Phi$.
  Wegen Satz \ref{thm:simplyTransitive}(c) existieren ein $\sigma \in \W$ und ein $\alpha \in \Delta$, sodass $\sigma(\alpha) = \gamma$ gilt.
  Aus Symmetriegründen nehmen wir an, dass $\alpha$ in $\Delta_1$ und somit $\gamma$ in $\Phi_1$ liegt.

  Es existiert eine Darstellung $\sigma = \sigma_{\alpha_1} \dots \sigma_{\alpha_t}$ ($\ast$) mit $\alpha_i \in \Delta$.
  Wir können ohne Beschränkung der Allgemeinheit $\alpha_1,\dots,\alpha_t \in \Delta_1$ annehmen.
  Denn, da nach Voraussetzung $(\Delta_1, \Delta_2) = 0$ gilt, folgt mit Lemma \ref{lem:orthogonalRoots} für die Darstellung ($\ast$), dass sich die Faktoren umordnen lassen zu
  \begin{displaymath}
    \sigma = \sigma_{\alpha_1} \dots \sigma_{\alpha_j} \sigma_{\alpha_{j+1}} \dots \sigma_{\alpha_t},
  \end{displaymath}
  wobei $1 \leq j < t$ und $\alpha_k \in \Delta_2$ für alle $k > j$ gilt.
  Da $\alpha$ aus $\Delta_1$ stammt und für $\alpha_k \in \Delta_2$ die Identität $\sigma_{\alpha_k}(\alpha) = \alpha - \langle \alpha, \alpha_k \rangle \alpha_k = \alpha$ gilt, folgt sogleich $\sigma = \sigma_{\alpha_1} \dots \sigma_{\alpha_j}$. 
  Damit ergibt sich schließlich $\gamma = \sigma(\alpha) = \sigma_{\alpha_1} \dots \sigma_{\alpha_j}(\alpha)$.

  Aus der Darstellung der Spiegelung an einer einfachen Wurzel folgt, dass $\gamma$ in $E_1 := \spn(\Delta_1)$ liegt.
  Es gilt also auch $\Phi_1 \subseteq E_1$.
  Mit der Bilinearität des Skalarproduktes folgt dann $(\Phi_1, \Phi_2) = 0$.
  Da das Skalarprodukt nicht ausgeartet ist, muss dann $\Phi_1$ oder $\Phi_2$ leer sein.
  Dies impliziert jedoch, dass im Widerspruch zur Voraussetzung $\Delta_1$ oder $\Delta_2$ leer ist.
\end{proof}

Das nächste Lemma beschäftigt sich mit bezüglich der Halbordnung $\succeq$ auf $\Phi$ maximalen Wurzeln. 
Ihre Existenz ist aufgrund der in \hyperref[it:R1]{(R1)} vorausgesetzten Endlichkeit von $\Phi$ immer gewährleistet. 

\begin{lem}
  \label{lem:maximalRoot}
  Sei $\Phi$ ein irreduzibles Wurzelsystem.
  Bezüglich der Halbordnung $\succeq$ auf $\Phi$ existiert genau eine maximale Wurzel $\beta$.
  Insbesondere folgt aus $\alpha \neq \beta$ auch $\hgt \alpha < \hgt \beta$ und $(\beta, \alpha) \geq 0$ für alle $\alpha \in \Delta$.
  In der Darstellung $\beta = \sum_{\alpha \in \Delta} k_\alpha \alpha$ sind alle Koeffizienten $k_\alpha$ strikt positiv.
\end{lem}

\begin{proof}
  Sei $\beta = \sum_{\alpha \in \Delta} k_\alpha \alpha$ eine maximale Wurzel bezüglich der Relation $\succeq$.
  Da jede einfache Wurzel bereits positiv ist, gilt zunächst auch $\beta \succeq 0$.
  Wir definieren $\Delta_1 = \{\alpha \in \Delta \mid k_\alpha > 0\}$ und $\Delta_2 = \{\alpha \in \Delta \mid k_\alpha = 0\}$.
  Aufgrund der Positivität von $\beta$ ist damit $\Delta = \Delta_1 \cup \Delta_2$ eine disjunkte Vereinigung.

  Dass $\Delta_1$ nicht leer ist, folgt bereits aus der Positivität von $\beta$.
  Wir wollen nun zeigen, dass $\Delta_2$ leer ist.
  Angenommen, $\Delta_2$ sei nicht leer.
  Es gilt dann $(\alpha', \alpha) \leq 0$ für alle $\alpha' \in \Delta_2$ und $\alpha \in \Delta$, denn sonst wäre nach Lemma \ref{lem:sumDiffRoot} auch $\alpha' - \alpha$ eine Wurzel, was jedoch \hyperref[it:B2]{(B2)} widerspricht.
  Damit folgt sogleich $(\alpha', \beta) = \sum_{\alpha \in \Delta} k_\alpha (\alpha', \alpha) \leq 0 $ ($\ast$).

  Da $\Phi$ nach Voraussetzung irreduzibel ist, muss ein $\alpha' \in \Delta_2$ existieren, welches nicht orthogonal zu $\Delta_1$ ist und für das $(\alpha', \alpha) < 0$ gilt.
  Ansonsten folgt nämlich aus der Zerlegung von $\Delta$ in $\Delta_1$ und $\Delta_2$ die Reduzibilität von $\Delta$, also nach Lemma \ref{lem:irreducibleRoot} auch die Reduzibilität von $\Phi$, was jedoch der Voraussetzung widerspricht.

  Seien also $\alpha \in \Delta_1$ und $\alpha' \in \Delta_2$ mit $(\alpha, \alpha') < 0$ gegeben.
  Damit folgt dann, dass in ($\ast$) sogar die strikte Ungleichung $(\alpha', \beta) < 0$ gilt.
  Lemma \ref{lem:sumDiffRoot} ergibt, dass $\alpha' +  \beta$ eine positive Wurzel ist.
  Dies widerspricht aber der vorausgesetzten Maximalität von $\beta$
  Die anfängliche Annahme $\Delta_2 \neq \emptyset$ muss also verworfen werden.
  Also sind alle Linearfaktoren $k_\alpha$ strikt positiv.

  Zudem zeigt der vorangehende Beweis, dass $(\alpha', \beta) \geq 0$ gilt für alle $\alpha' \in \Delta$ und dass zumindest ein $\alpha' \in \Delta$ existiert mit $(\alpha', \beta) > 0$.
    Denn $\Delta$ ist nach \hyperref[it:B2]{(B2)} eine Vektorraumbasis für $E$ und falls für alle $\alpha' \in \Delta$ die Gleichheit $(\alpha', \beta) = 0$ gilt, folgt $(v, \beta) = 0$ für alle $v \in E$ und damit $\beta = 0$, da Skalarprodukte nicht ausgeartet sind.

  Nun wollen wir die Eindeutigkeit der maximalen Wurzel zeigen.
  Sei dazu $\beta'$ eine weitere Wurzel.
  Nach \hyperref[it:B2]{(B2)} und den vorangehenden Ausführungen existiert wieder eine Darstellung der Form $\beta' = \sum_{\alpha \in \Delta} k_\alpha' \alpha$ mit $k_\alpha' > 0$ für alle $\alpha \in \Delta$.
  Ebenso existiert mindestens ein $\alpha' \in \Delta$ mit $(\alpha', \beta') > 0$.
  Damit folgt nun
  \begin{displaymath}
    (\beta, \beta') = \sum_{\alpha \in \Delta} k_\alpha (\alpha, \beta') \geq k_{\alpha'} (\alpha', \beta') > 0.
  \end{displaymath}
  Mit Lemma \ref{lem:sumDiffRoot} ergibt sich nun, dass auch $\beta - \beta'$ eine Wurzel ist, es sei denn, dass $\beta$ und $\beta'$ proportional sind. 
  In diesem Fall gilt $\beta = \beta'$, da beide Wurzeln voraussetzungsgemäß positiv sind.
  Angenommen, es gelten $\beta \neq \beta'$ und $\beta - \beta' \succeq 0$.
  Dann gilt auch $\beta \succeq \beta'$ im Widerspruch zur Maximalität von $\beta$.
  Analog zeigt man, dass auch die Annahme $\beta' \succeq \beta$ verworfen werden muss.
  Also gilt $\beta = \beta'$ und die maximale Wurzel ist eindeutig.
\end{proof}

\begin{lem}
  \label{lem:irreducibleGroupOp}
  Sei $\Phi$ ein irreduzibles Wurzelsystem.
  Dann ist auch die Gruppenoperation von $\W$ irreduzibel über $E$, es existieren also außer $\{0\}$ und $E$ keine weiteren $\W$\hyp{}invarianten Unterräume von $E$.
  Insbesondere erzeugt die $\W$\hyp{}Bahn einer beliebigen Wurzel $\alpha$ den gesamten Vektorraum $E$.
\end{lem}

\begin{proof}
  Sei $E'$ ein nichttrivialer $\W$\hyp{}invarianter Untervektorraum von $E$.
  Dann ist auch das orthogonale Komplement $(E')^\perp$ ein $\W$\hyp{}invarianter Unterraum.
  Ist nämlich $\alpha \in E'$ und $\beta \in (E')^\perp$, so gilt für $\sigma \in \W$ die Gleichung
  $\langle \alpha, \sigma(\beta) \rangle = \langle \sigma^{-1}(\alpha), \beta \rangle = 0$, da $\sigma^{-1}(\alpha) \in E'$ gilt.
  Es folgt also auch $\sigma(\beta) \in (E')^\perp$ und damit die $\W$\hyp{}Invarianz des Untervektorraums $(E')^\perp$. 
  Zudem gilt die Identität $E = E' \oplus (E')^\perp$.

  Dann ist aber auch $\Phi = (\Phi \cap E') \cup (\Phi \cap (E')^\perp$ eine Zerlegung von $\Phi$ in zueinander orthogonale Teilmengen.
  Da aber $\Phi$ als irreduzibles Wurzelsystem vorausgesetzt wurde muss eine der beteiligten Teilmengen leer sein.
  Angenommen, es gelte $\Phi \subseteq E'$. 
  Da Wurzelsysteme nach \hyperref[it:W1]{(W1)} den Vektorraum erzeugen, gilt $E = E'$. 
  Im Fall $\Phi \subseteq (E')^\perp$ folgt $E = (E')^\perp$.
  Damit haben wir den ersten Teil der Behauptung bewiesen.

  Es ist der Spann einer $\W$\hyp{}Bahn einer Wurzel $\alpha$ ein nicht trivialer $\W$\hyp{}invarianter Untervektorraum von $E$.
  Mit dem vorangehenden Beweis muss $\spn(\W\alpha) = E$ gelten.
  Damit ist die Behauptung vollständig bewiesen.
\end{proof}

\begin{bem}
  Der Beweis zeigt bereits auf, wie sich reduzible Wurzelsysteme aus irreduziblen Wurzelsystemen zusammensetzen.
  Zu jeder disjunkten Zerlegung des Wurzelsystems korrespondiert nämlich eine orthogonale Summe $\W$\hyp{}invarianter Untervektorräume.
  Eine Weiterführung dieses Gedankens zur Klassifikation von Wurzelsystemen findet sich in \cite[S.57ff.]{humphreys1972introduction}.
\end{bem}

\begin{lem}
  \label{lem:rootLength}
  Sei $\Phi$ irreduzibel.
  Es bezeichne $||\alpha|| := \sqrt{(\alpha,\alpha)}$ die vom Skalarprodukt auf $E$ induzierte Norm.
  Dann nimmt die Norm auf $\Phi$ höchstens zwei unterschiedliche Werte an.
  Zudem liegen alle Wurzeln mit gleicher Norm in einer gemeinsamen $\W$\hyp{}Bahn.
\end{lem}

\begin{proof}
  Seien $\alpha, \beta \in \Phi$.
  Dann können nicht alle Elemente aus der Bahn $\W\alpha$ orthogonal zu $\beta$ sein, da nach Lemma \ref{lem:irreducibleGroupOp} die Bahn $\W\alpha$ den Vektorraum $E$ erzeugt und $\beta$ sonst $0$ sein müsste.
  Aus \cite[S.45]{humphreys1972introduction} ist bekannt, dass  
  \begin{displaymath}
    \frac{(\beta,\beta)}{(\alpha,\alpha)} \in \left\{1,2,3,\frac{1}{2},\frac{1}{3}\right\} \tag{$\ast$}
  \end{displaymath}
  gilt, vorausgesetzt es ist $(\alpha, \beta) \neq 0$.
  Daraus folgt bereits der erste Teil der Behauptung.
  Denn, angenommen es existieren drei unterschiedliche Wurzellängen, dann müssten in ($\ast$) auch die Verhältnisse $\tfrac{2}{3}$ und $\tfrac{3}{2}$ auftauchen, was jedoch nicht der Fall ist.

  Seien nun $\alpha, \beta$ Wurzeln gleicher Norm.
  Da nach Lemma \ref{lem:irreducibleGroupOp} die Bahn $\W\alpha$ den Vektorraum $E$ erzeugt, ist $\beta$ nicht orthogonal zu allen Bahnelementen.
  Wir können daher ohne Beschränkung der Allgemeinheit wir $(\alpha, \beta) \neq 0$ annehmen.
  Ansonsten ersetzen wir $\alpha$ durch ein $\W$\hyp{}Bahnelement $\sigma(\alpha)$ mit $(\sigma(\alpha), \beta) \neq 0$.
  Im Falle $\alpha = \beta$ gibt es nichts zu zeigen, daher nehmen wir an, es gelte $\alpha \neq \beta$.
  Nun folgt unter Berücksichtigung von \cite[S.45]{humphreys1972introduction}, dass $\langle \alpha, \beta \rangle = \langle \beta, \alpha \rangle = \pm 1$ gelten muss.
  Aufgrund der im vorigen Ausdruck für $\alpha$ und $\beta$ geltenden Symmetrie ist der im Allgemeinen nur im ersten Argument lineare Term $\langle \alpha, \beta \rangle$ nun sogar bilinear.
  Wir können damit, falls nötig, $\beta$ auch durch $-\beta = \sigma_\beta(\beta)$ ersetzen, da diese Elemente in einer gemeinsamen $\W$\hyp{}Bahn liegen.
  Daher lässt sich annehmen, dass $\langle \alpha, \beta \rangle = 1$ gilt.
  Damit folgt
  \begin{align*}
    \sigma_\alpha \sigma_\beta \sigma_\alpha(\beta) 
    &= \sigma_\alpha \sigma_\beta(\beta - \langle \beta, \alpha \rangle \alpha) \\
    &= \sigma_\alpha \sigma_\beta(\beta - \alpha) \\
    &= \sigma_\alpha (-\beta - \alpha + \langle \alpha, \beta \rangle \beta) \\
    &= \sigma_\alpha (-\beta - \alpha + \beta) \\
    &= \sigma_\alpha (-\alpha) \\
    &= \alpha
  \end{align*}
  und es liegen $\alpha$ und $\beta$ in der selben $\W$\hyp{}Bahn.
\end{proof} 

\begin{bem}
  Ist das Wurzelsystem $\Phi$ irreduzibel und existieren Wurzeln unterschiedlicher Norm, so ist es auf Grundlage von Lemma \ref{lem:rootLength} möglich von \emph{langen} und \emph{kurzen} Wurzeln zu sprechen.
  Haben alle Wurzeln aus $\Phi$ dieselbe Norm, so ist es üblich alle Wurzeln als lang zu bezeichnen.
\end{bem}

\begin{lem}
  Sei $\Phi$ ein irreduzibles Wurzelsystem mit Wurzeln unterschiedlicher Norm.
  Dann ist die im Sinne von Lemma \ref{lem:maximalRoot} maximale Wurzel $\beta$ eine lange Wurzel.
\end{lem}

\begin{proof}
  Sei $\alpha \in \Phi$ gegeben.
  Da die Normabbildung eingeschränkt auf $\Phi$ nur zwei Werte annimmt, reicht es aus $(\beta, \beta) \geq (\alpha, \alpha)$ zu zeigen.
  Dazu ersetzen wir $\alpha$ und $\beta$ durch die nach \ref{lem:fundamentalDomain} eindeutig bestimmten Wurzeln aus $\W\alpha$ beziehungsweise $\W\beta$, die im Abschluss der Fundamentalkammer bezüglich $\Delta$ liegen. 
  Dies ist ohne Einschränkungen möglich, da Lemma \ref{lem:rootLength} garantiert, dass alle Bahnelemente dieselbe Norm besitzen.
  Nach Lemma \ref{lem:maximalRoot} gilt aufgrund der Maximalität von $\beta$ nun $\beta - \alpha \succeq 0$.
  Weiter ist $(\gamma, \beta - \alpha) \geq 0$ für alle $\gamma \in \overline{\mathfrak{C}(\Delta)}$, also insbesondere auch für $\gamma = \alpha$ und $\gamma = \beta$, was zu den Ungleichungen
  \begin{displaymath}
    (\beta, \beta) - (\beta, \alpha) \geq 0 \quad\text{und}\quad 
    (\beta, \alpha) - (\alpha, \alpha) 
    = (\alpha, \beta) - (\alpha, \alpha) \geq 0,
  \end{displaymath}
  also
  \begin{displaymath}
    (\beta, \beta)  \geq (\beta, \alpha) \quad\text{und}\quad 
    (\beta, \alpha) \geq (\alpha, \alpha)
  \end{displaymath}
  führt.
  Beide Ungleichungen zusammen ergeben folglich
  \begin{displaymath}
    (\beta, \beta) \geq (\beta, \alpha) \geq (\alpha, \alpha),
  \end{displaymath}
  was die Behauptung beweist.
\end{proof}

%\nocite{*}
\bibliographystyle{alphadin}
\bibliography{seminar16_fabian_gabel}
%\addcontentsline{toc}{chapter}{Literaturverzeichnis}
\addcontentsline{toc}{section}{Literaturverzeichnis}

      
\end{document}
